% !TEX root = ../test_multifile.tex
\section{Chapter One: Foundations}

This chapter covers the foundational material needed for the rest of the document.
We begin with a brief historical overview and then move into the core definitions
that will be referenced throughout subsequent chapters.

The concept of mathematical proof dates back to ancient Greece, where Euclid
formalized geometric reasoning in his Elements. Modern proof techniques build
upon this tradition with tools from set theory, logic, and abstract algebra.

Consider the following fundamental identity:
$$\sum_{i=1}^{n} i = \frac{n(n+1)}{2}$$

This result, sometimes attributed to Gauss, can be proved by induction. The base
case is trivial, and the inductive step follows from straightforward arithmetic.
We will use similar inductive arguments in later chapters when discussing
recursive data structures and algorithm correctness.

A second important concept is the notion of equivalence relations. A relation
$R$ on a set $S$ is an equivalence relation if it is reflexive, symmetric, and
transitive. Equivalence relations partition sets into disjoint classes, a fact
we will exploit heavily in Chapter Three.
